\documentclass[11pt,a4paper]{article}
\usepackage[margin=1in]{geometry}
\usepackage{amsmath,amssymb,amsthm}
\usepackage{algorithm}
\usepackage{algpseudocode}
\usepackage{booktabs}
\usepackage{hyperref}

\theoremstyle{definition}
\newtheorem{definition}{\textbf{Definition}}[section]
\newtheorem{theorem}{\textbf{Theorem}}[section]
\newtheorem{lemma}[theorem]{\textbf{Lemma}}
\newtheorem{remark}{\textbf{Remark}}[section]

\title{\textbf{Deep Q-Networks and Experience Replay (DQN)}}
\author{\textbf{Team Scaibu} \\ \small Email: \texttt{jaydeep.9664920749@gmail.com} \\ \small \textbf{Architecture and Core Engine Team}}
\date{\textbf{October 2026}}

\begin{document}
\maketitle

\section{Definitions and Cognitive Mental Models}

\subsection{Formal Mathematical Definition}
\textbf{Definition (Markov Decision Processes, Bellman Optimality, and Deep Q-Learning):}
Let a sequential decision task be formulated as a Markov Decision Process (MDP) defined by the 5-tuple $\mathcal{M} = (\mathcal{S}, \mathcal{A}, \mathcal{P}, \mathcal{R}, \gamma)$, where $\mathcal{S}$ is the state space, $\mathcal{A}$ is a discrete action space, $\mathcal{P}(s' | s, a)$ is the transition probability distribution, $\mathcal{R}(s, a)$ is the reward function, and $\gamma \in [0, 1)$ is the discount factor. The optimal action-value function $Q^*(s, a)$ represents the expected cumulative discounted return achievable by executing action $a$ from state $s$ and acting optimally thereafter:
{\boldmath
\begin{equation}
Q^*(s, a) = \max_\pi \mathbb{E}_\pi \left[ \sum_{t=0}^\infty \gamma^t R_{t+1} \;\middle|\; S_0 = s, A_0 = a \right]
\end{equation}
}
$Q^*(s, a)$ satisfies the fundamental \textbf{Bellman Optimality Equation}:
{\boldmath
\begin{equation}
Q^*(s, a) = \mathcal{R}(s, a) + \gamma \sum_{s' \in \mathcal{S}} \mathcal{P}(s' | s, a) \max_{a' \in \mathcal{A}} Q^*(s', a')
\end{equation}
}
In \textbf{Deep Q-Networks (DQN)}, $Q^*(s, a)$ is parameterized by a deep neural network $Q(s, a; \boldsymbol{\theta})$. To stabilize training on non-stationary, correlated continuous transition streams, DQN incorporates two foundational architectural mechanisms:
\begin{enumerate}
    \item \textbf{Experience Replay Buffer $\mathcal{D}$}: Transitions $e_t = (s_t, a_t, r_t, s_{t+1}, d_t)$ are stored in a cyclic memory buffer. Training samples are drawn uniformly at random, breaking temporal correlation.
    \item \textbf{Periodic Target Network $\boldsymbol{\theta}^-$}: To prevent policy oscillation from chasing a moving target, the regression target is computed using an older, frozen copy of the parameters $\boldsymbol{\theta}^-$:
    {\boldmath
    \begin{equation}
    y_t = r_t + (1 - d_t) \gamma \max_{a' \in \mathcal{A}} Q(s_{t+1}, a'; \boldsymbol{\theta}^-)
    \end{equation}
    }
\end{enumerate}
The network is trained via gradient descent minimizing the robust Huber loss on temporal difference (TD) error $\delta_t = y_t - Q(s_t, a_t; \boldsymbol{\theta})$:
{\boldmath
\begin{equation}
\mathcal{L}(\boldsymbol{\theta}) = \mathbb{E}_{(s, a, r, s', d) \sim \mathcal{D}} \left[ \ell_\delta\left( y - Q(s, a; \boldsymbol{\theta}) \right) \right]
\end{equation}
}
where $\ell_\delta(u) = \frac{1}{2} u^2$ for $|u| \le \delta$, and $\delta (|u| - \frac{1}{2} \delta)$ otherwise.

\subsection{Expanded Non-Technical Definition (Intuition for Anyone)}
\textbf{Intuitive Conceptual Definition:}
\begin{enumerate}
    \item \textbf{Learning by Trial and Error}: DQN is how a computer learns to master video games (like Atari Breakout or Pong) knowing nothing except screen pixels and the game score. It tries random button presses, observes when the score goes up, and learns which buttons lead to winning games.
    \item \textbf{The Fatal Trap of Consecutive Video Frames}: If an AI learns only from the immediate video feed, driving straight into a wall produces 50 frames in a row of crashing. The network overfits to crashing and forgets how to turn! Experience Replay acts like a photographic memory album: it shuffles thousands of past memories across hours of gameplay so the AI learns from a balanced, random variety of moments.
    \item \textbf{Freezing the Goalposts (Target Network)}: If you change your guess of what the future is worth every millisecond, you end up chasing your own tail. The Target Network freezes the evaluation rules for a few thousand steps, giving the student network a stable target to practice against.
    \item \textbf{Balancing Curiosity and Mastery ($\epsilon$-Greedy)}: With high probability the AI plays its best known move, but with small probability $\epsilon$ it tries a wild random action to discover hidden secrets.
\end{enumerate}

\subsection{Child-Friendly Mental Model (Physical Metaphor)}
Imagine a puppy learning an obstacle course:
\begin{itemize}
    \item \textbf{The Trial Run}: The puppy runs the course, sometimes getting a biscuit (reward $r$) and sometimes bumping into a cone.
    \item \textbf{The Dream Journal}: Every run is recorded on a trading card and tossed into a big jar (Replay Buffer).
    \item \textbf{Nighttime Training}: While the puppy sleeps, its trainer pulls out 32 random cards from the jar (some wins, some losses) and helps the puppy review: ``When you jumped here, you got a biscuit!''
    \item \textbf{The Master Dog}: By reviewing diverse memories rather than just the last 5 seconds, the puppy quickly becomes an agility champion!
\end{itemize}

\section{Core Mathematical Equations and Definitions}

\subsection{1-Step Bellman Optimality Target}
{\boldmath
\begin{equation}
y = 
\begin{cases}
r, & \text{if } s' \text{ is terminal} \\
r + \gamma \max_{a' \in \mathcal{A}} Q(s', a'; \boldsymbol{\theta}^-), & \text{otherwise}
\end{cases}
\end{equation}
}
\begin{itemize}
    \item \textbf{Formal Definition}: The bootstrapped 1-step Bellman optimality regression target evaluated using frozen target network parameters $\boldsymbol{\theta}^-$.
    \item \textbf{What It Means}: Estimates the true value of the current transition by summing the immediate reward with discounted future maximum return.
    \item \textbf{Why It Is Important}: Ground truth Q-values are unknown; the Bellman target serves as a self-consistent regression label for reinforcement learning.
\end{itemize}

\subsection{Temporal Difference Error Metric}
{\boldmath
\begin{equation}
\delta = y - Q(s, a; \boldsymbol{\theta})
\end{equation}
}
\begin{itemize}
    \item \textbf{Formal Definition}: The difference between the bootstrapped target return and the current online network prediction.
    \item \textbf{What It Means}: Measures the exact prediction surprise or error for having taken action $a$ in state $s$.
    \item \textbf{Why It Is Important}: Drives parameter updates via gradient descent: positive $\delta$ boosts action preference, while negative $\delta$ penalizes it.
\end{itemize}

\subsection{Robust Huber Loss Formulation}
{\boldmath
\begin{equation}
\ell_\delta(\delta) = 
\begin{cases}
\frac{1}{2} \delta^2, & \text{if } |\delta| \le \delta_{\text{th}} \\
\delta_{\text{th}} \left( |\delta| - \frac{1}{2} \delta_{\text{th}} \right), & \text{otherwise}
\end{cases}
\end{equation}
}
\begin{itemize}
    \item \textbf{Formal Definition}: The smooth combination of quadratic squared loss for small errors and linear absolute loss for large outliers.
    \item \textbf{What It Means}: Prevents gradient explosions when noisy environment rewards or rare events produce giant TD errors.
    \item \textbf{Why It Is Important}: Stabilizes deep Q-network backpropagation against extreme gradient swings during exploratory play.
\end{itemize}

\subsection{Greedy Action Selection Policy}
{\boldmath
\begin{equation}
a^* = \arg\max_{a \in \mathcal{A}} Q(s, a; \boldsymbol{\theta})
\end{equation}
}
\begin{itemize}
    \item \textbf{Formal Definition}: The deterministic policy selecting the action with maximal estimated expected return.
    \item \textbf{What It Means}: The optimal decision rule extracted from the learned value function during evaluation and exploitation.
    \item \textbf{Why It Is Important}: Converts learned scalar state-action values into physical control decisions.
\end{itemize}

\section{Rigorous Mathematical Analysis Checklist}

\begin{itemize}
    \item \textbf{Notation}: $Q(s, a)$ (action-value function), $\boldsymbol{\theta}$ (online parameters), $\boldsymbol{\theta}^-$ (target parameters), $\gamma$ (discount factor), $\mathcal{D}$ (replay buffer), $\delta$ (TD error).
    \item \textbf{Epistemic Status}: Seminal breakthrough by Mnih et al. (DeepMind, Nature 2015), the foundational milestone of deep reinforcement learning.
    \item \textbf{Dependency Graph}: $(s, a, r, s') \sim \mathcal{D} \to y \gets r + \gamma \max Q(s'; \boldsymbol{\theta}^-) \to \delta \gets y - Q(s, a; \boldsymbol{\theta}) \to \mathcal{L}_{\text{Huber}} \to \text{Adam step}$.
    \item \textbf{Dimensions}: Action space $|\mathcal{A}|$ (finite discrete), state space $\mathbb{R}^{d_s}$, scalar reward $r \in \mathbb{R}$.
    \item \textbf{Regimes}: Discrete-action control problems with dense or sparse reward feedback (Atari games, caching, discrete robotic routing).
    \item \textbf{Invariants}: Bellman target collapses to $y = r$ whenever $d_t = 1$ (terminal state).
    \item \textbf{Isomorphisms}: DQN is isomorphic to fitted Q-iteration where the non-linear regressor is an end-to-end convolutional/MLP neural network.
\end{itemize}

\section{Theoretical Theorems and Formal Proofs}

\begin{theorem}[Contraction Mapping of the Bellman Optimality Operator]
Let $\mathcal{B}: \mathbb{R}^{|\mathcal{S}| \times |\mathcal{A}|} \to \mathbb{R}^{|\mathcal{S}| \times |\mathcal{A}|}$ be the Bellman optimality operator defined by:
{\boldmath
\begin{equation}
(\mathcal{B} Q)(s, a) = \mathcal{R}(s, a) + \gamma \sum_{s' \in \mathcal{S}} \mathcal{P}(s' | s, a) \max_{a' \in \mathcal{A}} Q(s', a')
\end{equation}
}
Under the infinity norm $\|Q\|_\infty = \max_{s, a} |Q(s, a)|$, for discount factor $\gamma \in [0, 1)$, $\mathcal{B}$ is a $\gamma$-contraction mapping:
{\boldmath
\begin{equation}
\|\mathcal{B} Q_1 - \mathcal{B} Q_2\|_\infty \le \gamma \|Q_1 - Q_2\|_\infty
\end{equation}
}
Hence, by the Banach Fixed-Point Theorem, $\mathcal{B}$ has a unique fixed point $Q^*$ satisfying $\mathcal{B} Q^* = Q^*$.
\end{theorem}

\begin{proof}
Let $Q_1, Q_2 \in \mathbb{R}^{|\mathcal{S}| \times |\mathcal{A}|}$. For any state-action pair $(s, a)$:
{\boldmath
\begin{align}
|(\mathcal{B} Q_1)(s, a) - (\mathcal{B} Q_2)(s, a)| &= \left| \gamma \sum_{s'} \mathcal{P}(s' | s, a) \left( \max_{a'} Q_1(s', a') - \max_{a'} Q_2(s', a') \right) \right| \\
&\le \gamma \sum_{s'} \mathcal{P}(s' | s, a) \left| \max_{a'} Q_1(s', a') - \max_{a'} Q_2(s', a') \right|
\end{align}
}
Recall the standard real analysis identity: for any functions $f, g$, $|\max_x f(x) - \max_x g(x)| \le \max_x |f(x) - g(x)|$. Applying this mode-wise:
{\boldmath
\begin{equation}
\left| \max_{a'} Q_1(s', a') - \max_{a'} Q_2(s', a') \right| \le \max_{a'} |Q_1(s', a') - Q_2(s', a')| \le \|Q_1 - Q_2\|_\infty
\end{equation}
}
Substituting this back into the inequality:
{\boldmath
\begin{align}
|(\mathcal{B} Q_1)(s, a) - (\mathcal{B} Q_2)(s, a)| &\le \gamma \sum_{s'} \mathcal{P}(s' | s, a) \|Q_1 - Q_2\|_\infty \\
&= \gamma \|Q_1 - Q_2\|_\infty \sum_{s'} \mathcal{P}(s' | s, a) \\
&= \gamma \|Q_1 - Q_2\|_\infty
\end{align}
}
Since this inequality holds for all $(s, a) \in \mathcal{S} \times \mathcal{A}$:
{\boldmath
\begin{equation}
\|\mathcal{B} Q_1 - \mathcal{B} Q_2\|_\infty = \max_{s, a} |(\mathcal{B} Q_1)(s, a) - (\mathcal{B} Q_2)(s, a)| \le \gamma \|Q_1 - Q_2\|_\infty
\end{equation}
}
Since $\gamma < 1$, $\mathcal{B}$ is a strict contraction, guaranteeing convergence to a unique optimal value function $Q^*$.
\end{proof}

\begin{theorem}[Huber Loss Gradient Boundedness]
Let $\delta \in \mathbb{R}$ denote the TD error and let $\ell_\delta$ be the Huber loss with threshold $\delta_{\text{th}} > 0$. The gradient $\frac{\partial \ell_\delta}{\partial \delta}$ is strictly bounded:
{\boldmath
\begin{equation}
\left| \frac{\partial \ell_\delta}{\partial \delta} \right| \le \delta_{\text{th}}
\end{equation}
}
for all $\delta \in \mathbb{R}$. In contrast, standard MSE loss has unbounded gradient $|\frac{\partial \ell_{\text{MSE}}}{\partial \delta}| = |\delta| \to \infty$.
\end{theorem}

\begin{proof}
Differentiating the piecewise Huber loss with respect to $\delta$:
{\boldmath
\begin{equation}
\frac{\partial \ell_\delta}{\partial \delta} = 
\begin{cases}
\delta, & \text{if } |\delta| \le \delta_{\text{th}} \\
\delta_{\text{th}} \operatorname{sgn}(\delta), & \text{if } |\delta| > \delta_{\text{th}}
\end{cases}
\end{equation}
}
When $|\delta| \le \delta_{\text{th}}$, $|\frac{\partial \ell_\delta}{\partial \delta}| = |\delta| \le \delta_{\text{th}}$.
When $|\delta| > \delta_{\text{th}}$, $|\frac{\partial \ell_\delta}{\partial \delta}| = |\delta_{\text{th}} \operatorname{sgn}(\delta)| = \delta_{\text{th}}$.
Thus for all real $\delta \in \mathbb{R}$, $|\frac{\partial \ell_\delta}{\partial \delta}| \le \delta_{\text{th}}$, guaranteeing bounded gradient magnitudes and preventing gradient explosion during training.
\end{proof}

\section{Foundational Architectural Questions and Epistemic Meta-Questions}

\subsection{Architectural Question 1: Moving Target Instability}
\textbf{Question:} Why does training a deep Q-network without a separate target network $\boldsymbol{\theta}^-$ cause divergence?\\
\textbf{Answer:} If the online parameters $\boldsymbol{\theta}$ are used to compute both the prediction $Q(s, a; \boldsymbol{\theta})$ and the target $y = r + \gamma \max Q(s'; \boldsymbol{\theta})$, updating $\boldsymbol{\theta}$ simultaneously alters the ground truth labels. This creates a non-stationary feedback loop analogous to a dog chasing its own tail, leading to policy oscillation and divergent Q-value explosions.

\textbf{Meta-Question:} How frequently should the target network be synchronized?\\
\textbf{Answer:} Standard DQN updates $\boldsymbol{\theta}^- \gets \boldsymbol{\theta}$ periodically every $C = 1,000$ to $10,000$ environment steps (hard update). Modern implementations often use Polyak averaging (soft update): $\boldsymbol{\theta}^- \gets \tau \boldsymbol{\theta} + (1 - \tau) \boldsymbol{\theta}^-$ with $\tau \approx 0.005$ at every training step.

\subsection{Architectural Question 2: Overestimation Bias of the Maximum Operator}
\textbf{Question:} Why does standard DQN systematically overestimate action values?\\
\textbf{Answer:} The Bellman target uses the maximum operator: $\max_{a'} Q(s', a')$. Because neural network predictions contain stochastic estimation noise $\epsilon_{a'} \sim \mathcal{N}(0, \sigma^2)$, evaluating $\mathbb{E}[\max_{a'} (Q^* + \epsilon_{a'})] \ge \max_{a'} Q^*$. The expected value of the maximum of noisy estimates is strictly greater than or equal to the maximum of true values (Jensen's inequality), driving positive value bias that accumulates over time.

\textbf{Meta-Question:} How is overestimation bias resolved?\\
\textbf{Answer:} Through Double DQN (van Hasselt et al., 2015), which decouples action selection from action evaluation by using the online network to pick the action and the target network to evaluate its Q-value.

\section{Algorithmic Implementation Specification}

\begin{algorithm}[H]
\caption{Deep Q-Network (DQN) 1-Step Bellman Step and Huber Loss Evaluation}
\begin{algorithmic}[1]
\Require Current Q-values $\mathbf{q} \in \mathbb{R}^{|A|}$, action taken $a$, reward $r \in \mathbb{R}$, target next Q-values $\mathbf{q}_{\text{next}}^- \in \mathbb{R}^{|A|}$, discount $\gamma \in [0, 1]$, terminal flag $done$, Huber threshold $\delta_{\text{th}} > 0$
\Ensure Bellman target $y$, TD error $\delta$, Huber loss $\mathcal{L}$, greedy action $a^*$
\If{$done$ is true} \Comment{Terminal State Check}
    \State $y \gets r$
\Else
    \State $max\_q \gets \max_{0 \le a' < |A|} \mathbf{q}_{\text{next}}^-[a']$
    \State $y \gets r + \gamma \cdot max\_q$
\EndIf
\State $pred \gets \mathbf{q}[a]$
\State $\delta \gets y - pred$ \Comment{Compute TD Error}
\State $abs\_\delta \gets |\delta|$ \Comment{Evaluate Robust Huber Loss}
\If{$abs\_\delta \le \delta_{\text{th}}$}
    \State $\mathcal{L} \gets 0.5 \cdot \delta^2$
\Else
    \State $\mathcal{L} \gets \delta_{\text{th}} \cdot (abs\_\delta - 0.5 \cdot \delta_{\text{th}})$
\EndIf
\State $a^* \gets \arg\max_{0 \le i < |A|} \mathbf{q}[i]$ \Comment{Greedy Action Selection}
\State \Return $(y, \delta, \mathcal{L}, a^*)$
\end{algorithmic}
\end{algorithm}

\section{Big-$\Theta$ Complexity Analysis}

\begin{itemize}
    \item \textbf{Bellman Target Evaluation Time Complexity}: $\Theta(|A|)$ scalar operations to find the maximum action in $\mathbf{q}_{\text{next}}^-$.
    \item \textbf{TD Error and Huber Loss Complexity}: $\Theta(1)$ scalar operations.
    \item \textbf{Greedy Action Extraction}: $\Theta(|A|)$ operations.
    \item \textbf{Space Complexity}: $\Theta(1)$ working memory registers.
    \item \textbf{Parallel Depth}: $\mathcal{D} = \Theta(\log |A|)$.
\end{itemize}

\section{Single Canonical Repository Reference}

The complete deterministic scalar implementation, verification suite, and unit test benchmarks are published at:
\begin{center}
\url{https://github.com/Chief-Strategist-J/llm-obs-infra}
\end{center}

\end{document}
